% answers_short.tex —— 1977 年·湖北卷 · 简短答案速查
% 由 tools/make_short_answers.py 自 answers.tex 的 \gksolhead/\gkanswer 提取，
% 并与 1977/1977_数学真题_湖北_答案要点.md 逐题比对后生成，勿手改（改源 MD 后重跑）。

\begin{gknotice}
\textbf{使用说明}：本节只给最终答案便于快速核对，完整推导见解析版。答案为编写者独立推导，\textbf{非官方标准答案}。
\end{gknotice}

\gkpart{题目}
\begin{itemize}[leftmargin=2.6em]
  \item[\textbf{1.}] (1) $3$；$-27$；(2) $-2$；(3) $1$；(4) $x=2$；(5) $1-m$；(6) $-4<x\leqslant 2$ 或 $x<-4$（即 $x\leqslant 2$ 且 $x\ne-4$）；(7) $(x^2+y^2-xy)(x^2+y^2+xy)$；(8) 见解析版
  \item[\textbf{2.}] 与墙平行的一边长 $\dfrac{L}{2}$ 米、与墙垂直的两边各长 $\dfrac{L}{4}$ 米时，围成的矩形面积最大，最大面积为 $\dfrac{L^2}{8}$ 平方米。
  \item[\textbf{3.}] (1) $x-y+3=0$；(2) $2\sqrt2$
  \item[\textbf{4.}] $(x,y)=(3,9)$ 或 $(2,4)$。
  \item[\textbf{5.}] $y=2\sin\left(2x+\dfrac{\pi}{3}\right)$，周期为 $\pi$，振幅为 $2$。
\end{itemize}
